% GENERATED FILE. Edit canonical lessons, then run npm run generate:books.
% canonical-source-hash: fnv1a64:69d117c7de8868ea
% canonical-lessons: FA-C517-piyano, FA-C517-viyolon, FA-C517-tar, FA-C517-santur, FA-C517-tombak

\chapter{Musical Instruments}
\label{ch:fa-musical-instruments}
\hlchaptermodality{517}

\begin{chapteropening}
\textbf{By the end of this chapter:} \emph{I can name some musical instruments.}

Complete the last lesson of chapter 517: I can name some musical instruments.
\end{chapteropening}

\section[\texorpdfstring{\fa{پیانو} (piyâno)}{پیانو (piyâno)}]{\fa{پیانو} (piyâno) — a piano}

\label{lesson:FA-C517-piyano}



\begin{quote}
Before the new one: say the Persian for energetic, then the Persian for to be lazy.
\end{quote}

\subsection*{You'll want to know: \fa{پیانو}}
\textbf{\fa{پیانو}} (\emph{piyâno}) — \textquotedblleft{}a piano\textquotedblright{}.

\begin{quote}
\fa{دخترم} \fa{از} \fa{هفت} \fa{سالگی} \fa{پیانو} \fa{میزند}.
\end{quote}

\begin{quote}
\emph{dokhtaram az haft sâlegi piyâno mizanad.}
\end{quote}

\begin{quote}
My daughter has played the piano since she was seven.
\end{quote}

Persian strikes an instrument: \emph{piyâno zadan} 'to play the piano'.

\begin{cousinweb}
\emph{piyâno} is a modern loan, usually traced to French or Italian \emph{piano}. EtymDB has no entry for the Persian word.
\end{cousinweb}

\begin{grammarlens}[title={What you've built}]
One more word for telling, explaining and describing at B1.
\end{grammarlens}

\subsection*{Your turn}
\begin{itemize}
\raggedright
  \item \emph{Say it:} \emph{piyâno}
  \item \emph{Say it:} \emph{piyâno}, once more
  \item \emph{Say it:} a sentence of your own with \emph{piyâno}
  \item \emph{Recall:} say the Persian for energetic, then the Persian for to be lazy, then say \emph{piyâno} again
\end{itemize}

\begin{tcolorbox}[breakable,colback=teal!4,colframe=teal!35!black,title={Before you move on}]
What does \fa{پیانو} mean? (\textquotedblleft{}A piano\textquotedblright{}.) Say it once more.
\end{tcolorbox}
\section[\texorpdfstring{\fa{ویولن} (viyolon)}{ویولن (viyolon)}]{\fa{ویولن} (viyolon) — a violin}

\label{lesson:FA-C517-viyolon}



\begin{quote}
Before the new one: say the Persian for to be lazy, then the Persian for a piano.
\end{quote}

\subsection*{You'll want to know: \fa{ویولن}}
\textbf{\fa{ویولن}} (\emph{viyolon}) — \textquotedblleft{}a violin\textquotedblright{}.

\begin{quote}
\fa{او} \fa{در} \fa{کنسرت} \fa{ویولن} \fa{زد}.
\end{quote}

\begin{quote}
\emph{u dar konsert viyolon zad.}
\end{quote}

\begin{quote}
She played the violin at the concert.
\end{quote}

\emph{viyolon zadan} — to play the violin.

\begin{cousinweb}
\emph{viyolon} was borrowed from French \emph{violon}, 'violin' (EtymDB).
\end{cousinweb}

\begin{grammarlens}[title={What you've built}]
One more word for telling, explaining and describing at B1.
\end{grammarlens}

\subsection*{Your turn}
\begin{itemize}
\raggedright
  \item \emph{Say it:} \emph{viyolon}
  \item \emph{Say it:} \emph{viyolon}, once more
  \item \emph{Say it:} a sentence of your own with \emph{viyolon}
  \item \emph{Recall:} say the Persian for to be lazy, then the Persian for a piano, then say \emph{viyolon} again
\end{itemize}

\begin{tcolorbox}[breakable,colback=teal!4,colframe=teal!35!black,title={Before you move on}]
What does \fa{ویولن} mean? (\textquotedblleft{}A violin\textquotedblright{}.) Say it once more.
\end{tcolorbox}
\section[\texorpdfstring{\fa{تار} (târ)}{تار (târ)}]{\fa{تار} (târ) — a tar (a long-necked lute); a string}

\label{lesson:FA-C517-tar}



\begin{quote}
Before the new one: say the Persian for a piano, then the Persian for a violin.
\end{quote}

\subsection*{You'll want to know: \fa{تار}}
\textbf{\fa{تار}} (\emph{târ}) — \textquotedblleft{}a tar (a long-necked lute); a string\textquotedblright{}.

\begin{quote}
\fa{پدربزرگم} \fa{تار} \fa{میزد}.
\end{quote}

\begin{quote}
\emph{pedarbozorgam târ mizad.}
\end{quote}

\begin{quote}
My grandfather used to play the tar.
\end{quote}

\emph{târ} also means 'string, thread': the instrument is 'the string'.

\begin{cousinweb}
\emph{târ} is a native Persian word for a string or thread, and the instrument takes its name from it. EtymDB lists it without a source, so its older origin is not settled.
\end{cousinweb}

\begin{grammarlens}[title={What you've built}]
One more word for telling, explaining and describing at B1.
\end{grammarlens}

\subsection*{Your turn}
\begin{itemize}
\raggedright
  \item \emph{Say it:} \emph{târ}
  \item \emph{Say it:} \emph{târ}, once more
  \item \emph{Say it:} a sentence of your own with \emph{târ}
  \item \emph{Recall:} say the Persian for a piano, then the Persian for a violin, then say \emph{târ} again
\end{itemize}

\begin{tcolorbox}[breakable,colback=teal!4,colframe=teal!35!black,title={Before you move on}]
What does \fa{تار} mean? (\textquotedblleft{}A tar (a long-necked lute); a string\textquotedblright{}.) Say it once more.
\end{tcolorbox}
\section[\texorpdfstring{\fa{سنتور} (santur)}{سنتور (santur)}]{\fa{سنتور} (santur) — a santur (a hammered dulcimer)}

\label{lesson:FA-C517-santur}



\begin{quote}
Before the new one: say the Persian for a violin, then the Persian for a tar.
\end{quote}

\subsection*{You'll want to know: \fa{سنتور}}
\textbf{\fa{سنتور}} (\emph{santur}) — \textquotedblleft{}a santur (a hammered dulcimer)\textquotedblright{}.

\begin{quote}
\fa{صدای} \fa{سنتور} \fa{خیلی} \fa{آرام} \fa{است}.
\end{quote}

\begin{quote}
\emph{sedâ-ye santur kheyli ârâm ast.}
\end{quote}

\begin{quote}
The sound of the santur is very soft.
\end{quote}

Played with two light hammers: \emph{santur zadan}.

\begin{cousinweb}
\emph{santur} is the name of a Persian instrument, but EtymDB lists it without a source, and its origin is not settled.
\end{cousinweb}

\begin{grammarlens}[title={What you've built}]
One more word for telling, explaining and describing at B1.
\end{grammarlens}

\subsection*{Your turn}
\begin{itemize}
\raggedright
  \item \emph{Say it:} \emph{santur}
  \item \emph{Say it:} \emph{santur}, once more
  \item \emph{Say it:} a sentence of your own with \emph{santur}
  \item \emph{Recall:} say the Persian for a violin, then the Persian for a tar, then say \emph{santur} again
\end{itemize}

\begin{tcolorbox}[breakable,colback=teal!4,colframe=teal!35!black,title={Before you move on}]
What does \fa{سنتور} mean? (\textquotedblleft{}A santur (a hammered dulcimer)\textquotedblright{}.) Say it once more.
\end{tcolorbox}
\section[\texorpdfstring{\fa{تنبک} (tombak)}{تنبک (tombak)}]{\fa{تنبک} (tombak) — a tombak (a goblet drum)}

\label{lesson:FA-C517-tombak}



\begin{quote}
Before the new one: say the Persian for a tar, then the Persian for a santur.
\end{quote}

\subsection*{You'll want to know: \fa{تنبک}}
\textbf{\fa{تنبک}} (\emph{tombak}) — \textquotedblleft{}a tombak (a goblet drum)\textquotedblright{}.

\begin{quote}
\fa{او} \fa{در} \fa{گروه} \fa{موسیقی} \fa{تنبک} \fa{میزند}.
\end{quote}

\begin{quote}
\emph{u dar goruh-e musiqi tombak mizanad.}
\end{quote}

\begin{quote}
He plays the tombak in a music group.
\end{quote}

Spelled with \emph{n} but said with \emph{m} before \emph{b}: \emph{tombak}. Also \emph{tonbak}.

\begin{cousinweb}
\emph{tombak} is inherited: EtymDB traces it to a Middle Persian word for 'drum' (EtymDB).
\end{cousinweb}

\begin{grammarlens}[title={What you've built}]
One more word for telling, explaining and describing at B1.
\end{grammarlens}

\subsection*{Your turn}
\begin{itemize}
\raggedright
  \item \emph{Say it:} \emph{tombak}
  \item \emph{Say it:} \emph{tombak}, once more
  \item \emph{Say it:} a sentence of your own with \emph{tombak}
  \item \emph{Recall:} say the Persian for a tar, then the Persian for a santur, then say \emph{tombak} again
\end{itemize}

\begin{tcolorbox}[breakable,colback=teal!4,colframe=teal!35!black,title={Before you move on}]
What does \fa{تنبک} mean? (\textquotedblleft{}A tombak (a goblet drum)\textquotedblright{}.) Say it once more.
\end{tcolorbox}
